\chapter{自然数,有限集}

\section{自然数的定义}

\begin{definition}{有限基数,自然数,有限集,有限元素族}{}
	\begin{enumerate}
		\item 设$\mathfrak{a}$是基数.若$\mathfrak{a}\neq\mathfrak{a}+1$,则称$\mathfrak{a}$是有限基数或自然数.
		\item 设$E$是集合.若$\left|E\right|$是有限基数,则称$E$是有限集,$\left|E\right|$是$E$中元素的个数.
		\item 设$\left(x_{i}\right)_{i\in I}$是集合$E$的元素族.若$I$是有限集,则称$\left(x_{i}\right)_{i\in I}$是$E$中的有限元素族》
	\end{enumerate}
\end{definition}

\begin{remark}{术语说明}{}
	当我们说一系列对象有$m$个时(其中$m$是自然数),我们是指这些对象是一个基数为$m$的有限集的元素.基数为$m$的元素也被称为$m$元集.
\end{remark}

\begin{proposition}{}{}
	若$\mathfrak{a}$是基数,则$\mathfrak{a}$是有限基数$\iff$ $\mathfrak{a}+1$是有限基数.
\end{proposition}

\begin{remark}{}{}
	由于$0\neq 1$,因此$0$是自然数.进一步,$1$和$2$都是自然数.类似的,$2+1$和$\left(2+1\right)+1$都是自然数,分别记作$3$和$4$.
\end{remark}



















